% ==============================================================================================================
% ================================================ INPUT =======================================================
% ==============================================================================================================


% ================================================ COVER =======================================================

% Logos located at Resources/logos/...
\def\coversubjectlogo{AST-h-black.png} 
\def\coversubjectlogoscale{0.6} 
\def\largetitlename{Laboratorio 1: Análisis de Espectros Estelares}

% Select wether you want to include ORCID iDs and links
\newif\iforcidbool
% \orcidboolfalse   % or 
\orcidbooltrue

% Select citation and reference style:
%   \usenumcitetrue  -> numbered style [1], [2], ...
%   \usenumcitefalse -> default AASTeX author-year style (Author et al. YYYY)
%
% NOTE: after switching modes, delete .aux and .bbl before recompiling,
%       then run: pdflatex -> bibtex -> pdflatex -> pdflatex
\newif\ifusenumcite
% \usenumcitetrue   % or
\usenumcitefalse

% --- Authors ---
\def\coverauthors{
    {0009-0005-4685-695X}{Juan José Sánchez Medina},
    {}{Henrietta Leavitt},
    {}{Stephen Hawking}
}

% --- Professors ---
\def\coverprofessor{
    {0000-0000-0000-0000}{Isaac Newton}
}

% --- Assistants ---
\def\coverassist{
    {}{Marie Curie}
}

% --- Lab Date ---
\def\docinitdate{6 de noviembre}
% --- Final Date ---
\def\docfinalmonth{noviembre}

% --- Subject Details ---
\def\subjectcode{AST067}
\def\subjecttext{Laboratorio de Astronomía}


% ================================================ HEADER ======================================================
\shorttitle{Informe Laboratorio 01}
\shortauthors{Grupo 2070}

%% Include dates for submitted, revised, and accepted.
%%\received{February 1, 2025}
%%\revised{March 1, 2025}
%%\accepted{\today}
%%
%% Indicate AAS Journal the manuscript was submitted to.
%%\submitjournal{PSJ}
%% Note that this command adds "Submitted to " the argument.
%%
%% You can add a light gray and diagonal water-mark to the first page 
%% with this command:
%% \watermark{text}
%% where "text", e.g. DRAFT, is the text to appear.  If the text is 
%% long you can control the water-mark size with:
%% \setwatermarkfontsize{dimension}
%% where dimension is any recognized LaTeX dimension, e.g. pt, in, etc.
%%%%%%%%%%%%%%%%%%%%%%%%%%%%%%%%%%%%%%%%%%%%%%%%%%%%%%%%%%%%%%%%%%%%%%%%%%%%%%%%
%%
%% Use this command to indicate a subdirectory where figures are located.
%%\graphicspath{{./}{figures/}}

% =============================================== PREAMBLE =======================================================

% --- Custom Imports ---
\usepackage{amsmath}
\usepackage{cancel}
\usepackage{lipsum}
\usepackage{multirow}
\usepackage{soul}
\usepackage[colorinlistoftodos, spanish]{todonotes}
\usepackage{siunitx}

% --- Custom Commands ---
\newcommand{\vdag}{(v)^\dagger}
\newcommand\aastex{AAS\TeX}
\newcommand\latex{La\TeX}
\newcommand{\diff}{\mathop{}\!\mathrm{d}}
\newcommand{\cint}{%
  \mathop{%
    \vcenter{\hbox{\(\displaystyle\int\mkern-18mu C\)}}%
  }\nolimits%
}

\newcommand{\todoerror}[2][]{\todo[color=red!70, inline, #1]{#2}}
\newcommand{\todowarn}[2][]{\todo[color=orange, inline, #1]{#2}}
\newcommand{\todoremark}[2][]{\todo[color=yellow!80!black, inline, #1]{#2}}
\newcommand{\todohint}[2][]{\todo[color=green!20, inline, #1]{#2}}
\newcommand{\tododone}[2][]{\todo[color=blue!50, inline, #1]{#2}}

\makeatletter
\let\frontmatter@title@above=\relax
\makeatother